% TLDR: Place VeriBench and broader judge methodology in related work, not this paper's result claims.
\section{Related Work}
\label{sec:related}
\paragraph{VeriBench and formal specification coverage.}
VeriBench evaluates end-to-end coding and formal-verification agents, including whether their formal statements express required properties~\citep{veribench2026}.
Its tasks, reference artifacts and earlier human-rating study provide this paper's application context.
The broader solver benchmark, proof-success comparisons and operational coverage scorer remain separate contributions.
This paper extracts the human-rating judge question and preserves each historical scorer's input conditions, target and analysis.
Predicting expert coverage ratings is also different from forecasting acceptance by a bounded proof-search procedure: a search failure does not establish non-entailment.

\paragraph{Language-model judges.}
Human-agreement studies evaluate language-model judges for conversational preferences and report biases associated with position, verbosity and self-preference~\citep{zheng2023judging}.
Our task instead concerns formal statements, assumptions and reference coverage.
Agreement statistics remain useful, but an ordinal probability forecast additionally makes disagreement and distributional error explicit.
We compare training objectives with a common validation criterion rather than assume that any objective is empirically superior.

\paragraph{Scoring, selection and dependent observations.}
Proper scoring rules motivate probability forecasts whose expected score is optimized by the target distribution~\citep{gneiting2007strictly}.
Repeated model selection can bias subsequent performance assessment~\citep{cawley2010overfitting}, motivating separate validation and sealed test use.
Crossed-array resampling addresses dependence induced by shared identities~\citep{owen2012bootstrapping}; its suitability still depends on the actual assignment and target population.
These results guide the protocol but do not supply a universal number of tasks, solvers or raters.
